\section{Introduction}

Data attribution and data selection are often evaluated with the same
rank-and-retrain pipeline, but they answer different questions. Attribution
asks how an objective changes under a local perturbation of one training
example. Selection asks which budget-feasible data should shape an entire
training trajectory. A score evaluated once at uniform weights can therefore
be a useful direction without being a solution to the selection problem.

We make this distinction explicit through bilevel optimization over continuous
example weights. In this coordinate system, the value-function construction of
F2SA \citep{kwon2023f2sa} produces the difference between losses under two
lower-level solutions. Idealized RHO-LOSS \citep{mindermann2022rholoss}
produces the same per-example direction at a fixed outer state. Practical
RHO-LOSS is farther away: it replaces the updating joint solution with a
frozen holdout model, scores a batch-local candidate pool, takes a local quota,
and immediately updates a nonconverged target. These changes can be useful,
but their effects should not be attributed to the loss-difference formula
without a controlled decomposition.

The empirical question is thus not simply whether a value-function score beats
RHO-LOSS. It is whether optimizing the same continuous outer state beyond the
uniform point produces better fixed-budget data, and which parts of the
practical online protocol explain any remaining advantage. We answer this with
a fixed-horizon block design. After a shared warm start, every target runs for
five equal blocks. A selector may update at the first one through five block
boundaries, but all branches train for the same total epochs and examples. The
design directly compares one-shot selection, repeated scores reset to uniform,
and persistent projected optimization. It therefore distinguishes outer
memory from the simpler benefit of rescoring a changed model.

Our contributions are:
\begin{enumerate}
    \item We derive, in one native notation, the exact correspondence between
    the F2SA value-function direction and idealized RHO-LOSS, while marking the
    assumptions broken by the practical frozen-model estimator.
    \item We formulate one-shot attribution and multi-round selection as
    $R=1$ and $R=2,\ldots,5$ instances of one fixed-budget, fixed-horizon
    protocol, with reset and persistent outer states and matched random
    controls.
    \item We separate RHO's score approximation from its training protocol by
    adding global scoring, batch-local quota, immediate update, and online
    rescoring one at a time, then measure how the online batch scale changes
    accuracy, optimizer steps, coverage, and computation.
\end{enumerate}
The paper's main claim is deliberately conditional: persistent multi-round
optimization is supported only if it improves over both its one-shot branch
and the corresponding reset branch at equal target-training horizon. If a
batch-faithful RHO variant wins instead, the decomposition determines whether
the source is the frozen comparator, local diversity constraint, online
adaptation, or optimizer-step granularity.
